\documentclass[14pt,letterpaper,fleqn]{article}

% Packages
\usepackage[utf8]{inputenc}
\usepackage[T1]{fontenc}
\usepackage[english]{babel}
\usepackage[fleqn]{amsmath}
\usepackage{amssymb}
\usepackage{graphicx}
\usepackage{booktabs}
\usepackage{array}
\usepackage{longtable}
\usepackage{multirow}
\usepackage{hyperref}
\usepackage{geometry}
\usepackage{setspace}
\usepackage{natbib}
\usepackage{caption}
\usepackage{subcaption}
\usepackage{xcolor}
\usepackage{colortbl}
\usepackage{nomencl}
\usepackage{url}
\usepackage{adjustbox}

\geometry{
  top=2.5cm,
  bottom=2.5cm,
  left=3cm,
  right=3cm
}

\setlength{\mathindent}{0pt}

\onehalfspacing

% ----------------------------------------------------------------
\begin{document}

% ----------------------------------------------------------------
\begin{center}
  {\Large\bfseries Evaluating the efficiency and water recovery of direct contact
  membrane distillation (DCMD) systems with dual heat recovery and
  recirculation of retentate and permeate}

  \vspace{1em}

  Emerson B.\ dos Anjos$^{1,*}$,
  Abdul O.\ Cárdenas Gómez$^{1}$,
  Luz Elena P.\ Chenche$^{1}$,
  João A.\ Lima$^{1,2}$,
  Carolina P.\ Naveira-Cotta$^{1}$,
  Renato M.\ Cotta$^{1,3}$,
  Kleber M.\ Lisboa$^{4}$

  \vspace{0.8em}

  \small
  $^{1}$\textit{Laboratory of Nano \& Microfluidics and Microsystems -- LabMEMS,
    Mechanical Engineering Department, POLI \& COPPE/UFRJ,
    Federal University of Rio de Janeiro, Rio de Janeiro/RJ, Brazil}\\[0.3em]
  $^{2}$\textit{Renewable Energy Engineering Department, DEER/CEAR,
    Federal University of Paraíba, João Pessoa/PB, Brazil}\\[0.3em]
  $^{3}$\textit{General Directorate of Nuclear and Technological Development,
    DGDNTM, Brazilian Navy, Rio de Janeiro/RJ, Brazil}\\[0.3em]
  $^{4}$\textit{Laboratory of Thermal Sciences (LATERMO),
    Department of Mechanical Engineering -- TEM/UFF,
    Universidade Federal Fluminense, Niterói/RJ, Brazil}\\[0.5em]
  $^{*}$Corresponding author: \href{mailto:emersonbanjos@coppe.ufrj.br}{emersonbanjos@coppe.ufrj.br}
\end{center}

% ----------------------------------------------------------------
\begin{abstract}
A transient model for hollow-fiber direct contact membrane distillation (DCMD) modules
is employed to evaluate a desalination plant in six modes of operation, namely batch,
semi-batch, volume-control, salinity-control, continuous, and steady state. The main
goal is to reveal DCMD desalination configurations that present the best compromise
between maximum energy efficiency and minimum use of brackish water from wells, mainly
focusing on applications in sites under hydrological stress, like in the semi-arid
Brazilian northeastern region. The obtained results indicate that using an additional
heat exchanger for heat regeneration, combined with recirculation of the retentate
stream, enhances the system efficiency compared to a scenario that uses only the
distillate heat recovery. Additionally, the results also showed that if the main
interest is in obtaining high recovery ratios (high volume of produced distillate water
per volume of brackish water), the recirculating semi-batch and volume-controlling
operating modes are the best alternatives to produce distillated water. However, if
the interest is to minimize the thermal energy cost, the best option is the batch
operating mode.

\medskip
\noindent\textbf{Keywords:} Desalination; DCMD system; heat recovery; retentate recirculation;
energy efficiency.
\end{abstract}

\definecolor{headerblue}{RGB}{210, 225, 245}
\definecolor{sectiongray}{RGB}{235, 235, 235}

\newpage
\section*{Nomenclature}

\renewcommand{\arraystretch}{1.15}
\setlength{\tabcolsep}{6pt}
\setlength{\arrayrulewidth}{0.6pt}

% ---- PAGE 1: Latin Symbols + Greek Symbols ----
\resizebox{\textwidth}{!}{%
\begin{tabular}{|p{1.8cm}|p{5.5cm}||p{1.8cm}|p{5.5cm}|}
\hline
\multicolumn{2}{|l||}{\textbf{Latin Symbols (a--m)}} &
\multicolumn{2}{l|}{\textbf{Latin Symbols (n--z)}} \\
\hline
\textbf{Symbol} & \textbf{Description} &
\textbf{Symbol} & \textbf{Description} \\
\hline
$a$        & Specific surface area &
$N$        & Molar flux \\
\hline
$c_p$      & Specific heat at constant pressure &
$N_{fbr}$  & Number of fibers in the vessel \\
\hline
$C$        & Thermal capacity &
$Nu$       & Nusselt number \\
\hline
$C_s$      & Salt concentration &
$p$        & Pressure \\
\hline
$d$        & Diameter &
$Pr$       & Prandtl number \\
\hline
$D$        & Diffusivity / diffusion coefficient &
$q_c$      & Heat transferred to the cold sink \\
\hline
$f$        & Friction factor &
$q_h$      & Heat transferred from the hot source \\
\hline
$h$        & Convective heat transfer coefficient &
$q_{HX}$   & Heat transferred in the heat exchanger \\
\hline
$h_{fg}$   & Enthalpy / heat of vaporization &
$Q$        & Rate of heat transfer \\
\hline
$j_{wl}$   & Distillate flow rate &
$r$        & Radius \\
\hline
$k$        & Thermal conductivity &
$R$        & Universal gas constant \\
\hline
$L$        & Hollow-fiber module length &
$Re$       & Reynolds number \\
\hline
$m$        & Mass &
$T$        & Temperature \\
\hline
$\dot{m}$  & Mass flow rate &
$V$        & Volume \\
\hline
$M$        & Molar mass &
$\dot{V}$  & Volumetric flow rate \\
\hline
           &  &
$u$        & Stream velocity \\
\hline
           &  &
$z$        & Position along the module \\
\hline
\multicolumn{2}{|l||}{\textbf{Greek Symbols}} &
\multicolumn{2}{l|}{\textbf{Subscripts / Superscripts}} \\
\hline
\textbf{Symbol} & \textbf{Description} &
\textbf{Symbol} & \textbf{Description} \\
\hline
$\alpha$           & Membrane effectiveness (salt) &
$a$                & Relative to air \\
\hline
$\gamma$           & Relative to lumen or shell phase &
$amb$              & Relative to environment \\
\hline
$\delta$           & Membrane thickness &
$col$              & Relative to collector tank \\
\hline
$\varepsilon$      & Porosity &
$dist$             & Relative to distillate \\
\hline
$\varepsilon_{HX}$ & Heat exchanger efficacy &
$e$                & Relative to inlet \\
\hline
$\eta$             & Efficiency &
$ext$              & Relative to external \\
\hline
$\lambda$          & Rejection fraction, retentate stream &
$l$                & Relative to lumen \\
\hline
$\mu$              & Dynamic viscosity &
$m/memb$           & Relative to membrane \\
\hline
$\rho$             & Density &
$mix$              & Relative to mixture \\
\hline
$\tau$             & Tortuosity &
$o$                & Relative to outlet \\
\hline
$\phi$             & Salinity &
$s$                & Relative to shell \\
\hline
                   &  &
$vap$              & Relative to vapor \\
\hline
                   &  &
$w$                & Relative to water \\
\hline
                   &  &
$well$             & Relative to well \\
\hline
                   &  &
$\bar{(\cdot)}$    & Indicates averaged value \\
\hline
\end{tabular}
}

\newpage

% ---- PAGE 2: Acronyms ----
\section*{Nomenclature (continued)}

\renewcommand{\arraystretch}{1.15}
\setlength{\tabcolsep}{6pt}

\resizebox{\textwidth}{!}{%
\begin{tabular}{|p{1.8cm}|p{5.5cm}||p{1.8cm}|p{5.5cm}|}
\hline
\multicolumn{4}{|l|}{\textbf{Acronyms}} \\
\hline
\textbf{Symbol} & \textbf{Description} &
\textbf{Symbol} & \textbf{Description} \\
\hline
AGMD    & Air Gap Membrane Distillation &
PGMD    & Permeate Gap Membrane Distillation \\
\hline
CGMD    & Conductive Gap Membrane Distillation &
RO      & Reverse Osmosis \\
\hline
DCMD    & Direct Contact Membrane Distillation &
RR      & Recovery Ratio \\
\hline
DLR     & Distillate Limiting Ratio &
SEC     & Specific Energy Consumption \\
\hline
GOR     & Gained Output Ratio &
SGMD    & Sweeping Gas Membrane Distillation \\
\hline
FLR     & Feed Limiting Ratio &
V-AGMD  & Vacuum Air Gap Membrane Distillation \\
\hline
HICO    & Hot In, Cold Out &
VMD     & Vacuum Membrane Distillation \\
\hline
HOCI    & Hot Out, Cold In &
        &  \\
\hline
HX      & Heat Exchanger &
        &  \\
\hline
MD      & Membrane Distillation &
        &  \\
\hline
\end{tabular}
}

\clearpage
% ----------------------------------------------------------------
\section{Introduction}

Freshwater scarcity, mainly due to population growth, has attracted the attention of
researchers and industry for the development and improvement of techniques for the
treatment, purification, and desalination of salt water. Among the main desalination
techniques, those involving thermal separation methods, such as membrane distillation,
require more energy due to the need to vaporize the water since the water vaporization
enthalpy is about 2{,}400~kJ/kg \citep{Sharqawy2010}, while Gibbs free energy for
seawater, which governs the energy use of the reverse osmosis method, is 2.7~kJ/kg
\citep{Elimelech2011}. Therefore, one of the challenges that must be overcome for the
application of thermal methods such as membrane distillation (MD) in water desalination,
which aims to be competitive against the conventionally employed reverse osmosis (RO)
process in small and medium-sized systems ($<100$~m$^3$/day of water produced), is to
minimize the demand for thermal energy, improving thermal efficiency.

The membrane distillation process enables the use of residual heat from other thermal
processes or the heat available from renewable sources, such as solar collectors.
However, heat losses by conduction through the membrane further increase the thermal
consumption, which, when added to high-water consumption, limits the large-scale use of
these systems \citep{Khayet2011}. Nevertheless, according to \citet{Guan2015}, when
compared to conventional desalination processes, such as multi-stage flash distillation,
multi-effect distillation, and reverse osmosis, the MD process has several benefits:
low salinity sensitivity and high salt rejection (100\% content, theoretically), low
membrane vulnerability to fouling, greater performance under low severity operating
conditions, and the viability of using low-grade waste heat from other thermal processes
or renewable sources, such as solar collectors, as previously mentioned
\citep{Curcio2005, Charcosset2009, Drioli2015}.

Concerning its technology, according to the method of obtaining the vapor pressure
gradient through the membrane, the MD is classified in the following configurations
\citep{Zhang2015, Swaminathan2018}: DCMD (Direct Contact Membrane Distillation), AGMD
(Air Gap Membrane Distillation), SGMD (Sweeping Gas Membrane Distillation), VMD
(Vacuum Membrane Distillation), PGMD (Permeate Gap Membrane Distillation), CGMD
(Conductive Gap Membrane Distillation), V-AGMD (Vacuum-Enhanced Air Gap Membrane
Distillation), among other variants. Each configuration has its advantages and intrinsic
disadvantages, so the choice of one of them is still a research theme, which is usually
directed through metrics such as GOR (Gained Output Ratio) and SEC (Specific Energy
Consumption), which are parameters for energy efficiency analysis. The DCMD is the most
studied and straightforward configuration among them, widely used for desalination
studies and concentration of aqueous solutions.

In general, the increase in the thermal efficiency of MD systems is achieved by improving
the performance of membranes (by choosing membranes with low thermal conductivity,
adequate porosity, pore size and pore size distribution, and thickness), as well as by
reducing the amount of external heat to be supplied to the process
\citep{Gostoli1998, Khayet2006, AlObaidani2008, Yao2013, Li2014, Swaminathan2018}.

The heat supplied to the process can then be significantly reduced by the reuse of the
heat available in the output currents of the desalination module through external heat
exchangers, extending the applicability of the MD process and improving energy efficiency
\citep{Gilron2007, Lee2011, Hausmann2012, Lu2012, Zhang2015, GonzalezBravo2015,
Kim2015, Maheswari2015, Lin2014, Gustafson2018, Christie2020, Lisboa2021, Sampaio2022}.
Obviously, operating conditions such as temperature, the feed and distillate flow rates,
salinity of the feed stream, and its recirculation, among others, significantly affect
the MD process and can improve or reduce process efficiency
\citep{Khayet2000, Wang2009, Zhang2010, Fan2012, Alkhudhiri2012, Guan2014}.

Additionally, since MD presents a low recovery ratio per pass (ratio between the
production rate of pure water and the feed flow), using some form of recirculation of
the brine/retentate stream is mandatory to increase efficiency. In this way, by
considering a CGMD module, \citet{SwaminathanLienhard2018} describe four operation
strategies that enable high global recovery using a single-pass low recovery process
(batch recirculation, semi-batch recirculation, continuous recirculation, and continuous
multi-stage recirculation). The first two operations are classified as
transient/discontinuous processes, while the last two are steady state ones. According
to these authors, who studied the salinity variation from a lower limit of 70~g/kg to an
upper limit of 260~g/kg, the batch operation is the best option since GOR is about twice
for the same distillate flow when compared to the continuous recirculation operation.

\citet{Sampaio2022} modeled a seawater desalination DCMD plant with heat recovery of
the recirculating hot distillate stream by using a DCMD module composed by a bundle of
hollow-fiber tubes. He combined analytical solutions for the mass transfer problem
across the membranes with the finite volume method for solution the steady-state heat
transfer problem. Still, based on \citet{Christie2020} and \citet{Lin2014}, and to
extend their works, \citet{Sampaio2022} developed an approximation for the optimal
value of the ratio between permeate and feed flow rates, and for the maximum mass
recovery rate that, differently from the expressions developed by \citet{Christie2020}
and \citet{Lin2014}, take into consideration the membrane thermal efficiency.

Thus, based on the works of \citet{Swaminathan2018}, \citet{Christie2020},
\citet{Lisboa2021}, and \citet{Sampaio2022} and extending their analyses, the present
work aims at the development of physical-mathematical modeling, in the unsteady-state
regime, and subsequent numerical simulation of a hollow-fiber DCMD system with double
recirculation and heat recovery. The main goal is firstly to evaluate the potential of
heat recovery from the retentate stream to increase the energy efficiency of the system
to produce distilled water. Furthermore, an application of the model to study the
desalination process in regions with high seasonal hydrological stress, such as the
semi-arid region of the Brazilian northeastern. In addition to the classical batch,
semi-batch, continuous, and steady-state modes of operation, two intermediary modes,
namely, mixing tank volume controlling and mixing tank salt concentration controlling,
are also introduced. The aim is to verify their effect on the behavior of the system
efficiency and on the water recovery ratio (for minimizing the well brackish water
consumption). These controlling schemes can also be operated in conjunction with total or
partial recirculation, i.e., when the brine is rejected from the system with fraction
rejections equal or different to the unity, respectively.

Table~\ref{tab:litreview} presents, in a condensed way, a brief bibliographic review
of the main experimental and simulation works on the thermal efficiency of MD systems,
highlighting the main contributions of each of the works, especially, but not limited to,
those works that adopt heat recovery and recirculation of flow streams using hollow fiber
membrane configuration.

\begin{table}[htbp]
\centering
\caption{Literature review on thermal efficiency in MD systems by heat recovery.}
\label{tab:litreview}
\begin{adjustbox}{max width=\textwidth}
\begin{tabular}{lcccp{7cm}}
\toprule
\multirow{2}{*}{\textbf{Reference}} & \multicolumn{3}{c}{\textbf{Heat Recovery}} & \multirow{2}{*}{\textbf{Observations}} \\
\cmidrule{2-4}
 & \textbf{Distillate} & \textbf{Distillate} & \textbf{Retentate} & \\
\midrule
Fane \textit{et al.} (1987) & X & X & X & Experimental/Simulation -- Hollow fiber -- DCMD \\
Schneider \textit{et al.} (1988) & X & X & X & Experimental -- Hollow Fiber -- DCMD \\
Kurokawa and Sawa (1996) & X & X & X & Simulation -- Flat fiber -- DCMD \\
Li and Sirkar (2004) & --- & --- & --- & Experimental -- Hollow Fiber -- DCMD -- Crossflow \\
Gilron \textit{et al.} (2007) & X & X & X & Simulation -- Hollow fiber -- DCMD Cascade \\
Lee \textit{et al.} (2011) & X & X & X & Experimental/Simulation -- Hollow Fiber -- DCMD Cascade \\
Hausmann \textit{et al.} (2012) & --- & X & X & Experimental/Simulation -- Flat Fiber \\
He \textit{et al.} (2013) & X & X & X & Simulation -- Hollow fiber -- DCMD Cascade \\
Chung \textit{et al.} (2014) & X & X & X & Experimental/Simulation -- Hollow fiber -- DCMD cascade \\
Duong \textit{et al.} (2015) & --- & X & X & Experimental/Simulation -- Flat Fiber -- DCMD \\
Lin \textit{et al.} (2014) & X & X & --- & Simulation -- DCMD -- FLR/DLR \\
Guan \textit{et al.} (2015) & X & X & X & Simulation -- Hollow fiber -- DCMD \\
Zhang \textit{et al.} (2015) & X & X & X & Review on heat recovery in MD systems \\
Lokare \textit{et al.} (2018) & --- & X & X & Simulation -- DCMD -- Brine recycling \\
Christie \textit{et al.} (2020) & X & X & X & Simulation -- DCMD -- FLR/DLR \\
Lisboa \textit{et al.} (2021) & X & --- & --- & Simulation -- Hollow fiber -- DCMD \\
Sampaio (2022) & X & X & --- & Simulation -- Hollow fiber -- DCMD \\
\bottomrule
\end{tabular}
\end{adjustbox}
\end{table}

% ----------------------------------------------------------------
\section{System Description}

The MD desalination plant modelled in the present work is characterized by a water
desalination system that operates in the transient regime, containing a DCMD module (a
vessel with a bundle of hollow fiber membranes), two tanks (a mixing and collector), and
running with dual recirculation and dual heat recovery, where the heat available at the
output of the module from the feed and distillate streams are recovered through the heat
exchangers, HX-1 and HX-2, respectively, as shown in Figure~\ref{fig:system}. The heat
source and sink are constant-temperature heat reservoirs, so the feed and distillate
streams reach the specified values of the source/sink temperatures by supplying/extracting
exactly the corresponding amount of heat required.

\begin{figure}[htbp]
  \centering
  % \includegraphics[width=\linewidth]{fig1_system.png}
  \fbox{\textit{Figure 1 -- Design of the DCMD system with double heat recovery and
  partial/total recirculation of the retentate current.}}
  \caption{Design of the DCMD system with double heat recovery and partial/total
  recirculation of the retentate current.}
  \label{fig:system}
\end{figure}

At the initial time, the feed circuit is completely filled with brackish water at the
well salinity ($c_s^{well}$) and ambient temperature ($T_{amb}$). The distillate circuit
is filled with distilled water at the same ambient temperature. The feed stream leaves
the mixing tank at state~3 ($c_{s3}$, $V_3 = V_{feed}$, $T_3$) and enters HX-1, which
heats it to state~4 ($c_{s4}$, $V_4$, $T_4$) by receiving energy from the distillate
leaving the DCMD module at state~11 ($c_{s11}$, $V_{11}$, $T_{11}$). The mandatory use
of this heat exchanger (HX-1) greatly increases the efficiency of the MD system by
reducing the amount of heat ($q_h$) that must be supplied by the heater to raise the feed
stream temperature to the working temperature at state~5 ($c_{s5}$, $V_5$, $T_5 =
T_{feed}$).

The concentrated module effluent (the retentate) at state~6 ($c_{s6}$, $V_6$, $T_6$)
can be removed from the system (as a reject stream) or it can be recycled (recirculated).
In the present work, the computational code was built taking into consideration both
possibilities, simultaneously or not, depending on a flow rejection fraction ($\lambda$)
at state~9 ($c_{s9}$, $V_9$, $T_9$), which directly and dynamically impacts the
concentration of the solution in the tank and its volume, $V_{mix}$. When $\lambda = 1$,
the brine is completely rejected at state~9 (the MD system must be continuously fed from
the well on a permanent basis). Otherwise, when $\lambda < 1$, the retentate is partially
recycled/rejected.

When necessary, depending on salinity and volume of the mixing tank (or the operating
regime), brackish water can be pumped from the well at state~1 ($V_{well}$, $c_s^{well}$,
$T_{well}$), passing through HX-2 as a cold stream to gain energy carried by the effluent
retentate from the DCMD module at state~6, which will be cooled down to state~7 ($c_{s7}$,
$Q_7$, $T_7$). The stream at state~7 can be directed, either totally ($\lambda = 0$) or
partially ($0 < \lambda < 1$), to the mixing tank through state~8 ($c_{s8}$, $V_8$, $T_8$),
to be mixed with the heated brackish water stream coming from state~2.

To perform the present study, the operational ranges of the PLICA-7 DCMD module produced
by Aquastill for experiments in a relevant environment are adopted.

% ----------------------------------------------------------------
\section{Mathematical Model}

The implemented mathematical model incorporates heat and mass transfer in all system
components: mixing tank, distillate tank, DCMD module, heat source and sink, and heat
exchangers 1 and 2. For the DCMD module, the model proposed by \citet{Lisboa2021} was
extended to the unsteady-state regime. The authors treated the hollow-fiber DCMD module
as an equivalent porous medium composed of the following constituents: the porous shell
medium (where the feed stream flows), the porous lumen medium (where the distillate
stream flows), and the membrane (where the distillation process occurs).

An illustrative figure of the porous media model employed to analyze the heat and mass
transfer in the vessel is shown in Figure~\ref{fig:porousmedia}.

\begin{figure}[htbp]
  \centering
  % \includegraphics[width=\linewidth]{fig2_porous.png}
  \fbox{\textit{Figure 2 -- Illustration of the porous media models.}}
  \caption{Illustration of the porous media models. Cross sections of the: (a)~Porous
  medium, lumen side; (b)~Module and hollow-fiber membrane bundle; (c)~Porous medium,
  shell side. Figure~1 from \citet{Lisboa2021}. Reprinted with permission from Elsevier.
  Copyright 2021.}
  \label{fig:porousmedia}
\end{figure}

The hollow fiber membranes are considered to be uniformly distributed along a vessel
(module) cross section, with their axis parallel to the vessel axis. The module porosity
is defined as the ratio of the phase volume to the total vessel volume:
\begin{equation}
  \varepsilon_l = \frac{N_{fbr}\,\pi r_l^2}{\pi r_v^2}
  \label{eq:porl}
\end{equation}
\begin{equation}
  \varepsilon_s = \frac{\pi r_v^2 - N_{fbr}\,\pi r_s^2}{\pi r_v^2}
  \label{eq:pors}
\end{equation}

Here, $\varepsilon_\gamma$ is the porosity of phase $\gamma$ (with $\gamma = l$ for the
lumen phase and $\gamma = s$ for the shell phase), $N_{fbr}$ is the number of fibers in
the vessel, $r_\gamma$ is the radius of the $\gamma$ phase, and $r_v$ is the radius of
the vessel. Given that the membrane has a thickness $\delta$, the outer radius of the
membrane will be $r_s = r_l + \delta$.

\begin{figure}[htbp]
  \centering
  % \includegraphics[width=0.5\linewidth]{fig3_fiber.png}
  \fbox{\textit{Figure 3 -- Representation for the single-fiber model.}}
  \caption{Representation for the single-fiber model of the hollow-fiber membrane module.}
  \label{fig:fiber}
\end{figure}

The specific surface area, $a_\gamma$, used in heat and mass transfer calculations, is
the ratio of the surface area available for transport phenomena in each phase to the
vessel volume:
\begin{equation}
  a_\gamma = \frac{N_{fbr}\,2\pi r_\gamma}{\pi r_v^2}
  \label{eq:specsurf}
\end{equation}

Considering that the heat exchangers, the distillate and mixing tanks, and the DCMD
module are adiabatic, and that the salinity and temperature fields inside them are
uniform (except inside the DCMD module), the physical situation is mathematically
modelled as follows.

\subsection*{Mixture and Distillate Tanks}
\begin{align}
  \frac{dV_{mix}}{dt} &= \dot{V}_2 + \dot{V}_8 - \dot{V}_3 \label{eq:vmix}\\
  \frac{d(C_{s,mix}\,V_{mix})}{dt} &= \dot{V}_2\,C_{s2} + \dot{V}_8\,C_{s8} - \dot{V}_3\,C_{s3} \label{eq:csmix}\\
  \frac{d(T_{mix}\,V_{mix})}{dt} &= \dot{V}_2\,T_2 + \dot{V}_8\,T_8 - \dot{V}_3\,T_3 \label{eq:tmix}\\
  \frac{d(T_{col}\,V_{col})}{dt} &= \dot{V}_{12}\,T_{12} + \dot{V}_{jwl}\,T_{l,o} - \dot{V}_{11}\,T_{11} \label{eq:tcol}
\end{align}

\subsection*{DCMD Module}
\begin{align}
  \rho_s\,c_{p,s}\,\varepsilon_s\,\frac{\partial T_s}{\partial t} &=
    k_s\,\varepsilon_s\,\frac{\partial^2 T_s}{\partial z^2}
    - h_s\,a_s\,(T_s - T_{m,s}) - j_{wl}\,a_l\,h_{fg} \label{eq:modfeed}\\
  \rho_l\,c_{p,l}\,\varepsilon_l\,\frac{\partial T_l}{\partial t} &=
    k_l\,\varepsilon_l\,\frac{\partial^2 T_l}{\partial z^2}
    + h_l\,a_l\,(T_{m,l} - T_l) \label{eq:moddist}
\end{align}

with initial conditions:
\begin{equation}
  T_s(z,0) = T_l(z,0) = T_{amb} \label{eq:ic}
\end{equation}

and boundary conditions:
\begin{align}
  T_s(0,t) &= T_5\,,\quad T_l(L,t) = T_{col}(t) \label{eq:bc1}\\
  \frac{\partial T_s}{\partial z}\bigg|_{z=L} &= 0\,,\quad
  \frac{\partial T_l}{\partial z}\bigg|_{z=0} = 0 \label{eq:bc2}
\end{align}

In these expressions, $\rho_\gamma$ is the specific mass of phase $\gamma$, $c_{p,\gamma}$
is the specific heat at constant pressure, $k_\gamma$ is the thermal conductivity,
$T_\gamma$ is the bulk temperature, $T_{m,\gamma}$ is the interface temperature between
the hollow fiber membrane surface and the liquid phase $\gamma$, and $h_\gamma$ is the
convection heat transfer coefficient. The convective heat transfer coefficients for the
hot and cold streams are evaluated through the following empirical correlations:
\begin{equation}
  Nu_\gamma = \frac{(f_\gamma/8)\,Re_\gamma\,Pr_\gamma}
    {1.07 + 12.7\sqrt{f_\gamma/8}\,(Pr_\gamma^{2/3}-1)}
  \label{eq:Nu}
\end{equation}

where $Re_\gamma$ is the Reynolds number, $Pr_\gamma$ is the Prandtl number, and
$f_\gamma$ is the friction factor:
\begin{align}
  Re_\gamma &= \frac{\rho_\gamma\,u_\gamma\,d_{h,\gamma}}{\mu_\gamma}\,,\quad
  d_{h,l} = 2r_l\,,\quad d_{h,s} = 2(r_c - r_s) \label{eq:Re}\\
  f_l &= (0.790\ln Re_l - 1.64)^{-2}\,,\quad
  f_s = (0.790\ln Re_s - 1.64)^{-2} \label{eq:f}
\end{align}

The energy and mass balances for the remaining control volumes in the feed and distillate
streams provide the following relationships.

\noindent\textbf{For the temperatures:}
\begin{align}
  T_1 &= T_{well}\,,\quad T_3 = T_{mix} \label{eq:T13}\\
  T_6 &= T_s(L,t)\,,\quad T_{11} = T_l(0,t) \label{eq:T611}\\
  T_5 &= T_{feed}\,,\quad T_{12} = T_{col} \label{eq:T512}\\
  T_9 &= T_7 \label{eq:T9}\\
  T_{col,out} &= T_{col}\,,\quad T_{mix,out} = T_{mix} \label{eq:Tcol}
\end{align}

\noindent\textbf{For the flow rates:}
\begin{align}
  \dot{V}_1 &= \dot{V}_{well}\,,\quad \dot{V}_3 = \dot{V}_s \\
  \dot{V}_5 &= \dot{V}_3 = \dot{V}_6 = \dot{V}_s \\
  \dot{V}_9 &= \lambda\,\dot{V}_7\,,\quad \dot{V}_8 = (1-\lambda)\,\dot{V}_7 \\
  \dot{V}_{11} &= \dot{V}_{12} = \dot{V}_l
\end{align}

\noindent\textbf{For the salinities:}
\begin{align}
  C_{s1} &= C_{s,well}\,,\quad C_{s3} = C_{s,mix} \\
  C_{s6} &= C_{s7} = C_{s8} = C_{s9} \\
  C_{s,col} &= 0\,,\quad C_{s,dist} = 0
\end{align}

The temperatures $T_4$, $T_{12}$ and $T_2$, $T_7$ are obtained from the definitions
of the effectiveness of heat exchangers HX-1 and HX-2, respectively:
\begin{align}
  \varepsilon_{HX1} &= \frac{T_4 - T_3}{T_{11} - T_3} \label{eq:hx1}\\
  \varepsilon_{HX2} &= \frac{T_6 - T_7}{T_6 - T_1} \label{eq:hx2}\\
  T_2 &= T_1\,,\quad T_{12} = T_{col} \label{eq:T2T12}
\end{align}

In the salinity equations, $\alpha_{memb}$ is the effectiveness of the membrane in
retaining salt, taken as 100\% in the present study ($\alpha_{memb} = 1$). The properties
of the brackish well water are determined in each point of the system as a function of
temperature and salinity \citep{Sparrow2003}.

The temperatures of the membrane interfaces along the $z$-direction, $T_{m,l}(z,t)$
and $T_{m,s}(z,t)$, required in Eqs.~\eqref{eq:modfeed}--\eqref{eq:moddist}, are
evaluated iteratively through the following procedure:
\begin{align}
  Q_{m} &= h_s\,a_s\,(T_s - T_{m,s}) = Q_{cond} + Q_{lat} \label{eq:Qm}\\
  Q_{cond} &= \frac{k_m}{\delta}\,a_l\,(T_{m,s} - T_{m,l}) \label{eq:Qcond}\\
  Q_{lat} &= j_{wl}\,a_l\,h_{fg} \label{eq:Qlat}\\
  Q_{conv,l} &= h_l\,a_l\,(T_{m,l} - T_l) \label{eq:Qconvl}\\
  N_{wl} &= \frac{j_{wl}}{M_w} \label{eq:Nwl}
\end{align}

The density of the mixture at each position or control volume of the circuit, $\rho_i$,
is related to the salinity, $\phi_i$, through the density of pure water and salt at
$T_{amb} = 25^\circ$C ($\rho_{salt} = 2165$~kg/m$^3$ and $\rho_{H_2O} = 997.4$~kg/m$^3$):
\begin{equation}
  \rho_i = \rho_{H_2O} + \phi_i\,(\rho_{salt} - \rho_{H_2O})
  \label{eq:rho}
\end{equation}

Other fluid properties, such as specific heat at constant pressure, thermal conductivity,
viscosity, and enthalpy of vaporization, are also evaluated as a function of salinity and
temperature, according to the correlations proposed by \citet{Sparrow2003}.

\subsection{Modeling Mass Transfer in the Membrane}

To partially close the system of equations, the mass transfer inside the DCMD module
must be modeled in order to determine the water vapor mass flux through the membrane,
$j_{wl}$. For this, as performed by \citet{Lisboa2021}, the Dusty Gas model
\citep{Lawson1997}, one of the most widely used models for predicting flow of distilled
water through membranes, was employed:
\begin{equation}
  N_w + \frac{p_a}{p_w}\,\frac{N_w\,D_{wk}}{D_{wa0}} = -\frac{D_{wk}}{R\,T_m}\,\frac{dp_w}{dz}
  \label{eq:dustygas}
\end{equation}

Air is normally considered at rest in the membrane pores, so that $N_a \approx 0$.
Applying Dalton's law, one obtains the following expressions for the molar, $N_{wl}$,
and mass, $j_{wl}$, fluxes of distilled water through the hollow fiber membranes (based
on $r_l$):
\begin{align}
  N_{wl} &= \frac{D_{eff}}{R\,T_m\,\delta}\,(p_s - p_l) \label{eq:Nwl2}\\
  j_{wl} &= M_w\,N_{wl} \label{eq:jwl}
\end{align}

The Knudsen diffusivity of water vapor, $D_{wk}$, the molecular diffusivity of water in
air, $D_{wa0}$, and the effective molecular diffusivity, $D_{eff}$, are evaluated using
the classical relationships:
\begin{align}
  D_{wk} &= \frac{d_p}{3}\sqrt{\frac{8R\,T_m}{\pi\,M_w}} \label{eq:Dwk}\\
  D_{wa0} &= 1.895 \times 10^{-5}\,T_m^{2.072}\,P^{-1} \label{eq:Dwa0}\\
  D_{eff} &= \frac{\varepsilon_m}{\tau_m}\,\frac{D_{wk}\,D_{wa0}}{D_{wk} + D_{wa0}} \label{eq:Deff}
\end{align}

The total pressure inside the membrane pores, $P$, is usually the atmospheric pressure.
The partial vapor pressures of the salt streams at the membrane interfaces, $p_s$ and
$p_l$, are evaluated from the saturated vapor pressures, $p_{v\gamma}$, obtained by
using the Antoine equation, corrected for the respective salt mole fractions
($x_{salt,\gamma}$) and activity coefficients ($\alpha_\gamma$):
\begin{align}
  p_s &= \alpha_s\,(1-x_{salt,s})\,p_{v,s}(T_{m,s}) \label{eq:ps}\\
  p_l &= \alpha_l\,(1-x_{salt,l})\,p_{v,l}(T_{m,l}) \label{eq:pl}
\end{align}

The most common models for the membrane thermal conductivity are based on geometrical
distributions of the solid and gaseous phases. \citet{GarciaPayo2004} recommend that,
for porous media with porosity greater than 0.6, the Maxwell~I thermal conductivity model
be used. These models are mathematically described in terms of the solid matrix, $k_s$,
and fluid/gaseous, $k_f$, conductivities, as:
\begin{align}
  k_m &= \varepsilon_m\,k_f + (1-\varepsilon_m)\,k_s \quad\text{(parallel model)} \label{eq:km_par}\\
  k_m &= \left[\frac{\varepsilon_m}{k_f} + \frac{1-\varepsilon_m}{k_s}\right]^{-1} \quad\text{(series model)} \label{eq:km_ser}\\
  k_m &= k_f\,\frac{k_s + 2k_f - 2\varepsilon_m(k_f - k_s)}{k_s + 2k_f + \varepsilon_m(k_f - k_s)}\,,\quad
    k_f = \varepsilon_m\,k_{gas} + (1-\varepsilon_m)\,k_{polymer}
    \quad\text{(Maxwell~I model)} \label{eq:km_max}
\end{align}

A significant number of tortuosity models can be found in the literature
\citep{Ghanbarian2013}. \citet{Kim2021} analyzed 11 models of tortuosity. The most
common ones are:

\noindent\textbf{Geometric (Euclidian) Models:}
\begin{align}
  \tau_m &= \frac{(2-\varepsilon_m)^2}{\varepsilon_m} \quad\text{Mackies and Meares (1955)} \label{eq:tau_mm}\\
  \tau_m &= 1 + \frac{1-\varepsilon_m}{\varepsilon_m} \ln\frac{1}{1-\varepsilon_m} \quad\text{Beeckman (1990)} \label{eq:tau_beeck}
\end{align}

\noindent\textbf{Fractal Models:}
\begin{align}
  \tau_m &= \varepsilon_m^{-0.5} \quad\text{Elias-Kohav \textit{et al.} (1991)} \label{eq:tau_elias}\\
  \tau_m &= \varepsilon_m^{-0.5} \quad\text{Bruggeman (1935)} \label{eq:tau_brugg}
\end{align}

According to \citet{Curcino2022}, an adequate selection must be supported on the
knowledge of the morphological and structural properties of the membrane used in the
experimental desalination module.

\subsection{Efficiency Metrics}

To evaluate the desalination system and verify the improvement in energy efficiency, the
thermal efficiency $\eta_{th}$, the gain-output ratio $GOR$, the specific thermal energy
consumption $SEC_{th}$, and the recovery ratio $RR$ are employed and analyzed in
different operational scenarios. Those quantities are defined as:
\begin{align}
  \eta_{th} &= 1 - \frac{q_c}{q_h}\,,\quad
  GOR = \frac{h_{fg}\,\dot{m}_{dist}}{q_h}\,,\quad
  SEC_{th} = \frac{q_h}{\dot{V}_{dist}}\,,\quad
  RR = \frac{V_{dist}}{V_{well}} \label{eq:metrics}
\end{align}

where:
\begin{align}
  q_h &= \dot{m}_5\,c_{p5}\,(T_5 - T_4) \label{eq:qh}\\
  q_c &= \dot{m}_{12}\,c_{p12}\,(T_{12} - T_{col}) \label{eq:qc}\\
  q_{HX1} &= \dot{m}_{11}\,c_{p11}\,(T_{11} - T_{12}) \label{eq:qHX1}\\
  V_{dist} &= \int_0^{t_{end}} \dot{V}_{jwl}\,dt\,,\quad
  V_{well} = \int_0^{t_{end}} \dot{V}_{well}\,dt \label{eq:Vdist}
\end{align}

The global $GOR$ and $SEC_{th}$ parameters are taken as the average values up to the
simulation end time.

A computational code to solve the coupled system of ordinary and partial differential
equations was developed in the Fortran~95 language. The MMOLCH subroutine
\citep{IMSL2022}, for solving hyperbolic partial differential equation systems through
the method of lines, was employed, with a tolerance error set as $10^{-5}$. The
non-linear equations for the interface temperatures were solved, within an error target
prescribed at $10^{-11}$, with the routine ZBREN \citep{IMSL2022}.

% ----------------------------------------------------------------
\section{The System Operating Modes}

To control the operation of the dynamic system and optimize the desalination system water
production and brackish water consumption, a general algorithm based on volume and
salinity of the mixture in the mixing reservoir is applied at each instant of time. The
developed computational code is very flexible, allowing simulations of different
desalination system operation scenarios: steady-state regime (no recirculation), batch
recirculation, semi-batch recirculation, continuous recirculation (dynamic rejection
fraction/recirculation), and partial recirculation operation mode with volume or
concentration control.

For the unsteady-state regimes, the system can start the operation with no brine
rejection (total recirculation), $\lambda = \lambda_0 = 0$, or with a pre-defined
rejection (partial recirculation), $\lambda_0 \neq 0$. For these regimes, the
desalination system operates in such a way that the mixing tank, initially with a volume
of brackish water of $V_{mix_0} = 100$~L, has its mixture consumed by the DCMD module,
as distillate is produced over time.

Additionally, to analyze the influence of the external heat exchangers on the system
performance it is sufficient to vary their effectiveness ($\varepsilon_{HX1}$ and
$\varepsilon_{HX2}$). By choosing $\varepsilon_{HX1} = 0$ and/or $\varepsilon_{HX2} = 0$,
the heat exchangers are not included in the analysis.

\subsection{Batch and Semi-Batch Operations}

In batch and semi-batch operation modes, over each process cycle, the distillate is
continuously removed from the system, and the remaining solution that leaves the MD
module, with higher salinity, is directed to the feed tank ($\lambda = 0$). In the batch
operation process, there is no water pumping from the well ($\dot{V}_{well} = 0$) between
recycling, so the volume of the solution in the tank decreases, and the salt concentration
increases over time. In the semi-batch operation cycle, the volume of the mixing tank is
kept constant, by continuously adding fresh feed water from the well to the system, at
the same proportion the distillate is removed ($\dot{V}_{well} = \dot{V}_{jwl}$).

In both operation processes, the permeate production rate, as well as the heat transfer
rate, will reduce as the feed stream of the MD module becomes increasingly saturated with
salts. After the salinity of the solution in the mixing tank reaches a limiting value,
the high-salinity brine is flushed-out ($\lambda$ is set to 1) until a minimum volume
allowed is reached. Then, the power supply refills the tank from the well ($\dot{V}_{well}
= \dot{V}_s$) until the minimum concentration is reached, and the cycle restarts.

\subsection{Continuous Operation (Steady-State Operation with Partial Brine Recirculation)}

In the continuous operation regime, a fixed operational brine salinity is specified at
the MD module exit. Therefore, after the initial transients have elapsed and the solution
that leaves the module is at the required salinity (reaching the steady-state regime), it
must be continuously bled-out of the system at a specific fraction rejection ($\lambda$)
that depends on the MD module performance. The make-up water from the well
($\dot{V}_{well}$) is simultaneously added to the tank at a proportional flow rate that
keeps constant the salinity of the stream at the MD module exit, which is related to that
fraction rejection.

The rejection fraction and the volumetric flow rate from the well is obtained according
to the following salt and mixture mass balances at the mixing tank:
\begin{align}
  \lambda &= 1 - \frac{\dot{V}_{jwl}\,C_{s,well}}{(\dot{V}_s - \dot{V}_{jwl})\,(C_{s6} - C_{s,well})} \label{eq:lambda}\\
  \dot{V}_{well} &= \dot{V}_{jwl} + \lambda\,(\dot{V}_s - \dot{V}_{jwl}) \label{eq:Vwell}
\end{align}

\subsection{Steady-State Operation with No Brine Recirculation}

To simulate a steady state condition with no recirculation, it is sufficient to set
$\lambda = \lambda_0 = \lambda_1 = 1$ and $\dot{V}_{well} = \dot{V}_s$ all the time
during the simulations. In this situation, all retentate leaving the DCMD module is
rejected and the system is continuously fed by the well at constant feed flow rate and
salinity.

\subsection{Volume or Concentration Controlling}

In the mixing tank volume controlling operation, a minimum and a maximum allowed mixing
tank volume is set as the operation limits. The salinity of the water in the mixing tank
($c_{s,mix}(t)$) is continuously increasing with the distillate production, and its
volume ($V_{mix}(t)$) decreases. When the minimum volume is reached ($V_{min}$), the
tank is filled with brackish water from the well up to the maximum volume ($V_{max}$)
previously chosen. The process then repeats until finishing the desalination operation.

In the mixing tank salinity controlling operation, a minimum and a maximum allowed
mixing tank salinity is set as the operation limits. The salinity of the water in the
mixing tank ($c_{s,mix}(t)$) is continuously increasing with the distillate production,
while its volume ($V_{mix}(t)$) decreases since there is no brackish water flow from the
well. At the instant the maximum salinity allowed is reached ($C_{s,max}$), the pump from
the well is turned on to feed the mixing tank until the salinity of the water contained
in it falls to the minimum predefined value.

% ----------------------------------------------------------------
\section{Results}

\subsection{Algorithm Verification}

Before analyzing typical well operation scenarios, the developed computational code is
verified by comparing its results with available literature data. Since no transient
results are available, these comparisons are made with data for the steady-state regime.

Table~\ref{tab:membranes} presents the physical and geometric characteristics of the
membranes and modules employed in the comparisons. The experimental results of
\citet{Wang2008}, \citet{Yang2011} and numerical results of \citet{Guan2014} and
\citet{Kim2014} are employed for situations with no heat exchangers.

\begin{table}[htbp]
\centering
\caption{Characteristics (input data) of the membranes and modules used for comparison.}
\label{tab:membranes}
\small
\begin{tabular}{llllllll}
\toprule
\textbf{Parameter} & \textbf{Symbol} & \textbf{Unit}
  & \textbf{Wang \textit{et al.} (2008)}
  & \textbf{Kim (2014)}
  & \textbf{Yang \textit{et al.} (2011)}
  & \textbf{Guan \textit{et al.} (2014)} \\
\midrule
 & & & Experimental & Numerical & Experimental & Numerical \\
\midrule
Membrane inner radius & $r_l$ & mm & 0.3 & 0.3 & 0.46 & 0.49 \\
Membrane thickness & $t$ & mm & 0.11 & 0.11 & 0.275 & 0.240 \\
Polymer thermal conductivity & $k_p$ & W/(m$\cdot$K) & 0.18$^a$ & 0.18 & 0.18$^a$ & 0.18$^a$ \\
Average pore diameter & $d_p$ & mm & 0.17$^b$ & 0.17 & 0.145/0.191 & 0.2$^b$ \\
Porosity & $\varepsilon_m$ & --- & 0.738 & 0.738 & 0.83 & 0.85$^b$ \\
Shell inlet temperature & $T_{e,s}$ & $^\circ$C & 30--79.3 & 30--79.3 & 35--70 & 30--80 \\
Lumen inlet temperature & $T_{e,l}$ & $^\circ$C & 17.5 & 17.5 & 25 & 23.5--30 \\
Total pore pressure & $P$ & kPa & 101.325$^a$ & 101.325$^a$ & 101.325$^a$ & 101.325$^a$ \\
Shell side flow rate & $\dot{V}_s$ & L/min & 1.648$^c$ & 1.521 & 2.5--3.0 & 1.14--3.0 \\
Lumen side flow rate & $\dot{V}_l$ & L/min & 0.407$^c$ & 0.407 & 0.4 & 0.046--0.4 \\
Number of fibers & $N_{fbr}$ & --- & 30$^a$ & 30 & 20 & 22 \\
Vessel radius & $r_v$ & mm & 3.241$^a$ & 3.232 & 4.75 & 7.5 \\
Vessel length & $L$ & m & 0.200 & 0.200 & 0.225 & 0.200 \\
Salt mass fraction shell & $w_{NaCl,s}$ & \% & 3.5$^d$ & 3.5 & 3.5 & 7.0 \\
Lumen salt mass fraction & $w_{NaCl,l}$ & \% & 0.0 & 0.0 & 0.0 & 0.0 \\
Cold/hot fluid arrangement & --- & --- & HOCI & HOCI & HOCI & HICO \\
Direction of flow & --- & --- & Co-current & Countercurrent & Countercurrent & Countercurrent \\
Membrane area & $A_{memb}$ & cm$^2$ & 150 & 152.7 & 220 & 180.9 \\
Packing density & $\varphi$ & \% & 48$^a$ & 48 & 50 & 20.8$^a$ \\
\bottomrule
\end{tabular}

\smallskip\footnotesize
$^a$ Obtained from, or employed the calculation procedure in, \citet{Kim2014}.
$^b$ Assumed in comparison with other works.
$^c$ Calculated as a function of average velocities.
$^d$ \citet{Lisboa2021} used $w_{NaCl,s} = 2.5\%$ in their simulations.
\end{table}

Table~\ref{tab:verify} presents a verification of the results obtained with the present
work together with the work of \citet{Lisboa2021} and a comparison with the experimental
results of \citet{Wang2008}.

\begin{table}[htbp]
\centering
\caption{Verification of the computational code with the numerical results of
\citet{Lisboa2021} and comparison with the experimental data of \citet{Wang2008}.
$T_{e,l} = 17.5^\circ$C.}
\label{tab:verify}
\small
\begin{tabular}{cccccccc}
\toprule
$T_{e,s}$ [$^\circ$C] & $\Delta T$ [$^\circ$C]
  & \multicolumn{2}{c}{$T_{o,s}$ [$^\circ$C]}
  & \multicolumn{2}{c}{$T_{o,l}$ [$^\circ$C]}
  & \multicolumn{2}{c}{$j_{wl}$ [kg/(m$^2$$\cdot$h)]} \\
\cmidrule(lr){3-4}\cmidrule(lr){5-6}\cmidrule(lr){7-8}
 & & Present & Wang (2008) & Present & Wang (2008) & Present & Wang (2008) \\
\midrule
40.0 & 22.5 & 38.5 & 39.0 & 23.4 & 24.5 & 5.5  & $6.2 \pm 0.4$ \\
49.5 & 23.0 & 47.2 & 48.0 & 26.5 & 27.8 & 10.1 & $11.5 \pm 0.7$ \\
59.8 & 42.3 & 56.2 & 57.4 & 31.7 & 31.3 & 20.0 & $19.9 \pm 1.0$ \\
70.2 & 52.7 & 65.1 & 67.2 & 37.5 & 36.2 & 32.8 & $30.8 \pm 1.0$ \\
79.3 & 61.8 & 72.6 & 76.0 & 43.8 & 39.0 & 47.9 & $41.5 \pm 1.2$ \\
\bottomrule
\end{tabular}

\smallskip\footnotesize
* Based on the outer radius of the hollow fiber. $T_{Wang\_et\_al.}(2008)$: $T \pm 0.1^\circ$C.
\end{table}

The results show that the present methodology is consistent with the analyzed physics, at
least in the steady-state regime, when compared to the experimental results of
\citet{Wang2008}. This table also shows that the results obtained with the methodology
in transient regime converge, as would be expected for long times, to the numerical
results of \citet{Lisboa2021}.

Another verification is illustrated in Figure~\ref{fig:distflux}, which shows a
comparison of the results obtained in the present work for the produced distillate flux
with the experimental results of \citet{Wang2008} and the numerical results of
\citet{Kim2014}.

\begin{figure}[htbp]
  \centering
  % \includegraphics[width=0.6\linewidth]{fig4_distflux.png}
  \fbox{\textit{Figure 4 -- Comparison between the result obtained with the present
  methodology and the experimental results of Wang et al.\ (2008) and numerical results
  of Kim (2014) for the distillate flux.}}
  \caption{Comparison between the result obtained with the present methodology and the
  experimental results of \citet{Wang2008} and numerical results of \citet{Kim2014}
  for the distillate flux.}
  \label{fig:distflux}
\end{figure}

To illustrate the spatial temperature distribution along the DCMD module,
Figure~\ref{fig:tempguan} brings a comparison between the results obtained with the
present formulation and the numerical results of \citet{Guan2014} for the steady-state
temperatures along the feed channel, $T_s$, of the feed/membrane interface, $T_{m,s}$,
of the membrane/distillate interface, $T_{m,l}$, and of the distillate channel, $T_l$.
The inlet temperatures and flow rates of the feed and distillate streams employed in the
simulation were: $T_{e,s} = 80^\circ$C, $T_{e,l} = 30^\circ$C and $\dot{V}_s =
1.1425$~L/min, $\dot{V}_l = 0.04635$~L/min, respectively.

\begin{figure}[htbp]
  \centering
  % \includegraphics[width=0.6\linewidth]{fig5_tempguan.png}
  \fbox{\textit{Figure 5 -- Temperature distribution along the axial position in the
  DCMD module. Comparison with Guan et al.\ (2014).}}
  \caption{Temperature distribution along the axial position in the DCMD module in the
  feed, interface, and distillate regions. Comparison with the numerical results of
  \citet{Guan2014}.}
  \label{fig:tempguan}
\end{figure}

\subsection{Simulations with the Brackish Water Well}

\subsubsection{Preliminary Analysis}

For simulation of a real desalination system fed by brackish water wells, data from
wells commonly found in the northeastern region of Brazil were used, whose temperatures,
salt concentrations, typical flow rates, and other properties can be found in
\citet{Veiga2017}. Table~\ref{tab:wells} provides the characteristic parameter values
of 444 wells located in the Ceará state of Brazil.

\begin{table}[htbp]
\centering
\caption{Characteristics of wells in the state of Ceará/Brazil.}
\label{tab:wells}
\begin{tabular}{lcc}
\toprule
\textbf{Limits} & \textbf{Salt concentration [g$_s$/L]} & \textbf{Flow rates [m$^3$/h]} \\
\midrule
Maximum & 13 & 15 \\
Minimum & 0.199 & 0.10 \\
Average & 2.532 & 3.45 \\
\textbf{Adopted} & \textbf{3.5} & \textbf{3} \\
\bottomrule
\end{tabular}

\smallskip\footnotesize $T_{well} = T_{amb} = 25^\circ$C
\end{table}

The parameters relating to the characteristics of the DCMD module and the membrane are
those provided in Table~\ref{tab:module}. The flow arrangement of the module is carried
out in the HOCI configuration (hot fluid in the shell and cold fluid in the lumen) in
counter-current.

\begin{table}[htbp]
\centering
\caption{Input data, membrane and module, used in the simulation of the wells.}
\label{tab:module}
\begin{tabular}{lllll}
\toprule
\textbf{Parameter} & \textbf{Symbol} & \textbf{Unit} & \textbf{Lisboa \textit{et al.} (2021)} \\
\midrule
Membrane inner radius & $r_l$ & mm & 0.3 \\
Membrane thickness & $t$ & mm & 0.1 \\
Polymer thermal conductivity & $k_p$ & W/(m$\cdot$K) & 0.2 \\
Average pore diameter & $d_p$ & mm & 0.2 \\
Membrane porosity & $\varepsilon_m$ & --- & 0.8 \\
Inlet temperature on shell side & $T_{e,s}$ & $^\circ$C & 80 \\
Inlet temperature on the lumen side & $T_{e,l}$ & $^\circ$C & 25 \\
Total pore pressure & $P$ & kPa & 101.325 \\
Shell-side flow rate & $\dot{V}_s$ & L/min & 30--50--80 \\
Flow rate on the lumen side & $\dot{V}_l$ & L/min & 30--50--80 \\
Number of fibers & $N_{fbr}$ & --- & 10{,}000 \\
Vessel radius & $r_v$ & mm & 50 \\
Vessel length & $L$ & m & 1 \\
Salt mass fraction on shell side & $w_{NaCl,s}$ & \% & 0.035 \\
Lumen-side salt mass fraction & $w_{NaCl,l}$ & \% & 0 \\
Membrane area & $A_{memb}$ & m$^2$ & 7.2 \\
Packing & $\varphi$ & \% & 0.64 \\
\bottomrule
\end{tabular}
\end{table}

To illustrate a typical dynamic behavior of the heat transfer process inside the DCMD
module, Figure~\ref{fig:tempprofiles} shows the bulk temperature profiles of the retentate
and distillate streams along the respective channels inside the module, at four instants
of time (approximately $2.8\times10^{-4}$~s, 3.4~s, 16~s, and 30~s), for the flow
condition $\dot{V}_l = \dot{V}_s = 50$~L/min.

\begin{figure}[htbp]
  \centering
  % \includegraphics[width=0.6\linewidth]{fig6_tempev.png}
  \fbox{\textit{Figure 6 -- Time evolution of temperature profiles inside the DCMD module.}}
  \caption{Time evolution of temperature profiles inside the DCMD module for $\dot{V}_l
  = \dot{V}_s = 50$~L/min.}
  \label{fig:tempprofiles}
\end{figure}

As can be seen from this figure, the time evolution of the heat transfer inside the DCMD
module is very fast, since the steady-state condition is reached in just a few seconds,
taking no more than 30~s for the module to attain this regime.

Figures~\ref{fig:tempsteady}a and \ref{fig:tempsteady}b illustrate bulk temperature
profiles of the feed and permeate channels along the module for different conditions of
volumetric flow rate of the streams, considering steady-state single-pass flow in the
DCMD module, and no heat recovery ($\lambda_0 = \lambda = 1$, $\dot{V}_{well} = \dot{V}_s$,
$\varepsilon_{HX1} = \varepsilon_{HX2} = 0$).

\begin{figure}[htbp]
  \centering
  % \includegraphics[width=\linewidth]{fig7_tempsteady.png}
  \fbox{\textit{Figure 7 -- Steady-state temperature profiles along the axial position
  in the module DCMD.}}
  \caption{Steady-state temperature profiles along the axial position in the module DCMD
  for: a)~different feed and distillate flow rates; b)~equal distillate and feed flow
  rates ($\dot{V}_l = \dot{V}_s$).}
  \label{fig:tempsteady}
\end{figure}

For an intuitive understanding of how the operational and configurational decisions affect
the energy efficiency, \citet{Christie2020}, based on \citet{Lin2014}, introduced a
critical mass flow rate ratio, $\alpha^*$, to define the theoretical condition of maximum
water recovery factor in steady-state single-pass DCMD systems, $R_{max}$:
\begin{align}
  \alpha^* &= \frac{\dot{m}_s\,c_{p,s}}{\dot{m}_l\,c_{p,l}}\bigg|_{\text{critical}}\,,\quad
  \alpha = \frac{\dot{m}_s\,c_{p,s}}{\dot{m}_l\,c_{p,l}} \label{eq:alpha}\\
  R_{max}^{FLR} &> R_{max}^{DLR} \label{eq:Rmax}
\end{align}

For $\alpha < \alpha^*$, the module operates in the DLR regime. Conversely, for
$\alpha > \alpha^*$, the module operates in the FLR regime.

Table~\ref{tab:table6} brings some results obtained for the averaged theoretical maximum
water recovery factor, distillate mass flux, temperature polarization coefficient, and
DCMD module thermal efficiency ($GOR$ and $SEC_{th}$). Three scenarios related to the
use of the external heat exchangers for two feed salinity levels and different feed and
distillate flow rates were studied.

\begin{table}[htbp]
\centering
\caption{Influence of the operational conditions on the averaged DCMD module efficiency
metrics. Steady-state regime (8~h).}
\label{tab:table6}
\scriptsize
\begin{tabular}{lcccccccccccc}
\toprule
\textbf{Flow Condition} & $\alpha$ & $\alpha^*$ & \textbf{Oper.}
  & $T_{thr}$ [$^\circ$C] & $Z_{thr}$ [m]
  & $RR$ & $\dot{V}_{jwl}$ & TPC & $\eta_{th}$
  & \multicolumn{3}{c}{$GOR$ / $SEC_{th}$ [kWh/L]} \\
\cmidrule(lr){11-13}
 & & & & & & & [L/min] & [-] & [\%]
  & $\varepsilon_{HX1,2}=0$
  & $\varepsilon_{HX1}=0.8$; $\varepsilon_{HX2}=0$
  & $\varepsilon_{HX1}=\varepsilon_{HX2}=0.8$ \\
\midrule
\multicolumn{13}{l}{\textit{$w_{NaCl,s} = 3.5\%$}} \\
$\dot{V}_l=30 < \dot{V}_s=80$ & 0.38 & 0.88 & DLR & 80.0 & 0.63 & 2.5 & 2.6 & 0.73 & 73 & 0.28 / 2.5 & 0.41 / 1.7 & 1.4 / 0.48 \\
$\dot{V}_l \approx 44 < \dot{V}_s = 50$ & \textbf{0.88} & \textbf{0.88} & \textbf{FLR} & --- & --- & \textbf{5.0} & \textbf{3.3} & \textbf{0.56} & \textbf{67} & \textbf{0.56 / 1.2} & \textbf{1.9 / 0.36} & \textbf{3.3 / 0.21} \\
$\dot{V}_l = \dot{V}_s = 50$ & 1.0 & 0.88 & FLR & --- & --- & 5.2 & 3.4 & 0.59 & 66 & 0.58 / 1.2 & 1.8 / 0.38 & 2.5 / 0.27 \\
\midrule
\multicolumn{13}{l}{\textit{$w_{NaCl,s} = 12.0\%$}} \\
$\dot{V}_l \approx 42 < \dot{V}_s = 50$ & \textbf{0.74} & \textbf{0.74} & \textbf{FLR} & --- & --- & \textbf{2.2} & \textbf{1.7} & \textbf{0.72} & \textbf{37} & \textbf{0.30 / 2.3} & \textbf{0.88 / 0.79} & \textbf{1.6 / 0.43} \\
$\dot{V}_l = \dot{V}_s = 80$ & 0.88 & 0.74 & FLR & --- & --- & 2.9 & 3.6 & 0.68 & 50 & 0.40 / 1.7 & 0.96 / 0.72 & 1.5 / 0.46 \\
\bottomrule
\end{tabular}

\smallskip\footnotesize
(a) Feed-side threshold temperature, $T_{thres,s}$. (b) Distillate-side threshold temperature, $T_{thres,l}$.
\end{table}

\subsubsection{Operating Regime Analysis}

After characterizing the DCMD desalination system, investigation of the behavior of the
system for different operating modes is now performed. In all situations analyzed, the
FLR configuration to which $\dot{V}_l = \dot{V}_s = 50$~L/min is adopted.

Figures~\ref{fig:batchgor1} and \ref{fig:batchgor2} illustrate the dynamic behavior of
the $GOR$ and $SEC_{th}$ metrics for the batch operation mode and for two different
maximum allowed salt concentration of the brine in the mixing tank analyzed, $C_{s,max}$
(120~g$_s$/L$_{mix}$ and 260~g$_s$/L$_{mix}$, respectively). The averaged values
indicated in these figures were obtained for an operating time of 8~hours.

\begin{figure}[htbp]
  \centering
  \fbox{\textit{Figure 8 -- Evolution of the instantaneous GOR and SECth for $C_{s,max} = 120$~g$_s$/L$_{mix}$.}}
  \caption{Evolution of the instantaneous a)~$GOR$ and b)~$SEC_{th}$, for
  $C_{s,max} = 120$~g$_s$/L$_{mix}$. Batch operating mode.}
  \label{fig:batchgor1}
\end{figure}

\begin{figure}[htbp]
  \centering
  \fbox{\textit{Figure 9 -- Evolution of the instantaneous GOR and SECth for $C_{s,max} = 260$~g$_s$/L$_{mix}$.}}
  \caption{Evolution of the instantaneous a)~$GOR$ and b)~$SEC_{th}$, for
  $C_{s,max} = 260$~g$_s$/L$_{mix}$. Batch operating mode.}
  \label{fig:batchgor2}
\end{figure}

Figure~\ref{fig:modes} illustrates the temporal evolution of the main DCMD system
parameters for the different desalination system operating modes (batch, semi-batch,
volume-controlling, continuous and steady-state): mixture salt concentration in the
mixing tank $C_{s,mix}$, volumetric flow rate of the produced distillate $\dot{V}_{jwl}$,
mixture volume in the mixing tank $V_{mix}$, mixture temperature in the mixing tank
$T_{mix}$, volumetric flow rate of the brackish water well $\dot{V}_{well}$, and
rejection fraction $\lambda$.

\begin{figure}[htbp]
  \centering
  \fbox{\textit{Figure 10 -- Dynamic evolution of the instantaneous desalination system parameters
  for different operating modes.}}
  \caption{Dynamic evolution of the instantaneous desalination system parameters for
  different operating modes and $C_{s,max} = 120$~g$_s$/L$_{mix}$: a)~$C_{s,mix}$,
  b)~$\dot{V}_{jwl}$, c)~$V_{mix}$, d)~$T_{mix}$, e)~$\lambda$, and f)~$\dot{V}_{well}$.}
  \label{fig:modes}
\end{figure}

According to these figures, the cycling time for the batch mode is about 1789~s ($\approx
30$~min) while for the V-control and semi-batch modes are 4800~s ($\approx 1$~h~20~min)
and 6090~s ($\approx 1$~h~41~min), respectively.

In an attempt to identify the most appropriate operating mode of the DCMD desalination
system, comparisons are performed in Figures~\ref{fig:sec} and \ref{fig:rr}, where the
averaged specific thermal energy cost, $SEC_{th}$, and water recovery ratio, $RR$, are
displayed, respectively.

\begin{figure}[htbp]
  \centering
  \fbox{\textit{Figure 11 -- Effect of the system operation mode on the averaged specific thermal energy cost.}}
  \caption{Effect of the system operation mode on the averaged specific thermal energy
  cost, $SEC_{th}$, for different mixture salinity.}
  \label{fig:sec}
\end{figure}

\begin{figure}[htbp]
  \centering
  \fbox{\textit{Figure 12 -- Effect of the system operation mode on the water recovery ratio.}}
  \caption{Effect of the system operation mode on the water recovery ratio, $RR$, for
  different mixture salinity.}
  \label{fig:rr}
\end{figure}

Overall, according to these figures, one must seek a balance between the specific energy
cost and the water recovery ratio, since there exist, for the same maximum operational
salinity, a trade-off between these metrics. When it comes to the energy cost, and for
the same limiting mixture salinity, the best operating mode is the batch, followed by
the semi-batch, steady-state, and continuous modes, as indicated in the literature
\citep{SwaminathanLienhard2018}. On the other hand, when it comes to water recovery,
Figure~\ref{fig:rr} reveals that the best choice is the semi-batch operating mode,
followed by the V-control and batch modes.

Figure~\ref{fig:voldist} shows the potential of each operating mode in producing distillate
water when the desalination system is operated for 8~h.

\begin{figure}[htbp]
  \centering
  \fbox{\textit{Figure 13 -- Effect of the system operation mode on the produced volume
  of distillate water.}}
  \caption{Effect of the system operation mode on the produced volume of distillate water,
  for different mixture salinity.}
  \label{fig:voldist}
\end{figure}

In general, the desalination system can produce up to 1{,}200~L of distillate water
running in any of the operating modes for the lower salinity level. On the other hand,
when the salinity increases, it impacts negatively, especially for the steady-states modes
(continuous and steady-state), when only about 750~L of distillate water is produced.

Therefore, if the desalination system must be implemented in areas with hydrologic
seasonal stress, as the semi-arid region of the Brazilian northeastern, the results point
out that the semi-batch could be the preferred choice for brackish water desalination,
since the main interest is in saving salt water from the well.

% ----------------------------------------------------------------
\section{Conclusions}

The present work developed a mathematical modeling and computer simulation of a DCMD
desalination system with double heat recovery from the distillate stream and the feed
stream, together with total or partial recirculation of the retentate. The main goal was
to employ the developed code to identify the best scenarios for operation of the system
in terms of maximizing its thermal efficiency at the same time minimizing the use of
brackish water from wells, seeking optimized production of the amount of distilled water.

In general, the following conclusions can be summarized from the present results:

\begin{itemize}
  \item Heat recovery from the distillate stream already produces a significant increase
    in the energy efficiency of the desalination system, being a mandatory condition for
    operation of DCMD-based systems.

  \item Use of an additional heat exchanger to recover heat from the retentate stream
    also improves the energy efficiency of the process, especially when the system
    operates in the FLR condition with mass flow rate near the critical mass flow rate.

  \item Peaks in $GOR$ values and lows in $SEC$ values can be observed when the well and
    an external heat exchanger coupled to its stream operate simultaneously, so double
    heat recovery and partial recirculation of the retentate play an important role in
    the competitiveness of the membrane distillation process compared to other forms of
    desalination, especially in long operation processes.

  \item The amount of heat supplied to the system is strongly impacted by the rejection
    fraction of the retentate stream.

  \item The system operating mode has a strong impact on the energy consumption and water
    recovery of the DCMD desalination system.

  \item The lowest specific thermal energy cost is associated with the batch operation
    mode of the desalination plant, however, in regions with hydrologic seasonal stress,
    the results point out that the semi-batch could be the preferred choice for brackish
    water saving. Therefore, its choice must be based on information as where the system
    has to be implanted and what are the availability of energy and brackish water
    salinity.
\end{itemize}

Extension to other desalination MD methods, like AGMD and V-AGMD, are currently in
development, in order to simulate the entire experimental TRL-6 desalination unit
installed in the LabMEMS/UFRJ.

% ----------------------------------------------------------------
\bibliographystyle{abbrvnat}
\bibliography{ref}

\end{document}
